%% ╔══════════════════════════════════════════════════════════════════╗
%% ║         CORIOTLAB — INFORME DE ACTIVIDADES                      ║
%% ║         Documento interno · flexible                            ║
%% ║         Dos variantes: PERIÓDICO y GUÍA                        ║
%% ╚══════════════════════════════════════════════════════════════════╝

%% ═══════════════════════════════════════════════════════════════════
%% ZONA DE CONFIGURACIÓN — editar aquí antes de compilar
%% ═══════════════════════════════════════════════════════════════════

%% VARIANTE ACTIVA: "periodico" o "guia"
\newcommand{\TipoVariante}{periodico}

%% Período (solo para variante periódico): Semanal | Quincenal | Mensual
\newcommand{\TipoPeriodo}{Semanal}

%% Campos comunes
\newcommand{\DocTitulo}{Sistema de Monitoreo IoT con ESP32}
\newcommand{\DocNumero}{IA-2025-007}
\newcommand{\DocAutor}{Santiago Leal}
\newcommand{\DocCargo}{Investigador --- CORIOTLAB}
\newcommand{\DocProyecto}{Proyecto de Monitoreo Ambiental}
\newcommand{\DocPeriodo}{16 al 20 de junio de 2025}
\newcommand{\DocFecha}{\today}
\newcommand{\DocVersion}{1.0}

%% ═══════════════════════════════════════════════════════════════════

\documentclass[12pt, letterpaper]{article}

%% --- Codificación y lenguaje
\usepackage[spanish, es-nodecimaldot, es-noshorthands]{babel}

%% --- Fuentes (XeLaTeX --- cargar por nombre de familia Windows)
\usepackage{fontspec}
\setmainfont{Inter}
\newfontfamily\FuenteTitulo{MuseoModerno}
\newfontfamily\FuenteCodigo{Space Mono}

%% --- Geometría
\usepackage[
  top=2.8cm, bottom=2.8cm,
  left=2.5cm, right=2.5cm,
  headheight=1.8cm
]{geometry}

%% --- Color
\usepackage[table]{xcolor}
\definecolor{AzulITM}     {HTML}{102D69}
\definecolor{Magenta}     {HTML}{C14894}
\definecolor{AzulDigital} {HTML}{56ACDE}
\definecolor{GrisPizarra} {HTML}{2F2F2F}
\definecolor{GrisClaro}   {HTML}{F2F2F2}
\definecolor{GrisMedio}   {HTML}{AAAAAA}
\definecolor{GrisLinea}   {HTML}{DDDDDD}
\definecolor{Blanco}      {HTML}{FFFFFF}
\definecolor{VerdeTarea}  {HTML}{1A5C35}

%% --- Gráficos y membrete
\usepackage{graphicx}
\usepackage{tikz}
\usepackage{eso-pic}

%% --- Encabezado y pie
\usepackage{fancyhdr}
\usepackage{lastpage}

%% --- Tablas
\usepackage{longtable}
\usepackage{booktabs}
\usepackage{array}
\usepackage{colortbl}
\usepackage{multirow}

%% --- Listas
\usepackage{enumitem}

%% --- Cajas
\usepackage[most]{tcolorbox}

%% --- Código
\usepackage{listings}

%% --- Hipervínculos
\usepackage[hidelinks]{hyperref}

%% --- Espaciado de párrafos
\usepackage{parskip}
\setlength{\parskip}{6pt}

%% -----------------------------------------------------------------
%% MEMBRETE DE FONDO (8% opacidad en todas las páginas)
%% -----------------------------------------------------------------
\AddToShipoutPictureBG{%
  \begin{tikzpicture}[remember picture, overlay]
    \node[opacity=0.08] at (current page.center) {%
      \includegraphics[width=\paperwidth, height=\paperheight,
        keepaspectratio=false]%
        {../../membretes/membrete_azul.png}%
    };
  \end{tikzpicture}%
}

%% -----------------------------------------------------------------
%% ENCABEZADO Y PIE DE PÁGINA
%% -----------------------------------------------------------------
\pagestyle{fancy}
\fancyhf{}

\fancyhead[L]{%
  {\FuenteTitulo\small\color{AzulITM}\textbf{CORIOTLAB}}%
  {\color{GrisMedio}\small\ \ |\ \ }%
  {\small\color{GrisPizarra}\DocTitulo}%
}
\fancyhead[R]{%
  {\small\color{GrisMedio}\DocNumero}%
}
\fancyfoot[L]{%
  {\footnotesize\color{GrisMedio}www.coriotlab.co\ \ |\ \ info@coriotlab.co}%
}
\fancyfoot[C]{%
  {\footnotesize\color{GrisMedio}\DocFecha}%
}
\fancyfoot[R]{%
  {\footnotesize\color{GrisMedio}Página \thepage\ de \pageref{LastPage}}%
}

\renewcommand{\footrulewidth}{0.4pt}
\renewcommand{\footrule}{%
  {\color{GrisLinea}\rule{\headwidth}{0.4pt}}%
  \vskip 4pt%
}
\renewcommand{\headrulewidth}{1.2pt}
\renewcommand{\headrule}{%
  \vspace{2pt}%
  {\color{AzulITM}\rule{\headwidth}{1.2pt}}%
}

%% -----------------------------------------------------------------
%% ESTILOS DE SECCIONES
%% -----------------------------------------------------------------
\usepackage{titlesec}
\titleformat{\section}
  {\FuenteTitulo\large\color{AzulITM}\bfseries}
  {\thesection.}{0.5em}{}[{\color{AzulITM}\titlerule[0.8pt]}]

\titleformat{\subsection}
  {\FuenteTitulo\normalsize\color{GrisPizarra}\bfseries}
  {\thesubsection.}{0.5em}{}

%% -----------------------------------------------------------------
%% MACRO: PORTADA
%% -----------------------------------------------------------------
\newcommand{\PortadaActividades}[2]{%
  \thispagestyle{empty}
  \begin{tikzpicture}[remember picture, overlay]
    \fill[AzulITM] (current page.north west)
      rectangle ([yshift=-4.5cm] current page.north east);
    \node[anchor=west, text=white] at
      ([xshift=2.5cm, yshift=-1.5cm] current page.north west)
      {\FuenteTitulo\large\textbf{#1}};
    \node[anchor=east, text=AzulDigital] at
      ([xshift=-2.5cm, yshift=-1.5cm] current page.north east)
      {\FuenteTitulo\Large\textbf{#2}};
    \draw[AzulDigital, line width=2pt]
      ([yshift=-4.5cm] current page.north west) --
      ([yshift=-4.5cm] current page.north east);
  \end{tikzpicture}%
  \vspace{3.8cm}

  \begin{center}
    {\FuenteTitulo\Huge\color{AzulITM}\textbf{\DocTitulo}}\\[0.6em]
    {\FuenteTitulo\large\color{GrisPizarra}\DocProyecto}\\[2em]
    {\color{GrisLinea}\rule{\linewidth}{0.6pt}}\\[1.2em]
    \begin{tabular}{ll}
      {\small\color{GrisMedio}Autor:}   & {\small\color{GrisPizarra}\DocAutor} \\[3pt]
      {\small\color{GrisMedio}Cargo:}   & {\small\color{GrisPizarra}\DocCargo} \\[3pt]
      {\small\color{GrisMedio}Período:} & {\small\color{GrisPizarra}\DocPeriodo} \\[3pt]
      {\small\color{GrisMedio}Versión:} & {\small\color{GrisPizarra}\DocVersion} \\[3pt]
      {\small\color{GrisMedio}Fecha:}   & {\small\color{GrisPizarra}\DocFecha} \\
    \end{tabular}
  \end{center}
  \newpage
}

%% -----------------------------------------------------------------
%% MACRO: ESTADO DE TAREA
%% -----------------------------------------------------------------
\newcommand{\EstadoCompletado}{%
  \textcolor{VerdeTarea}{\small\textbf{Completado}}}
\newcommand{\EstadoEnCurso}{%
  \textcolor{AzulITM}{\small\textbf{En curso}}}
\newcommand{\EstadoPendiente}{%
  \textcolor{GrisMedio}{\small\textbf{Pendiente}}}
\newcommand{\EstadoBloqueado}{%
  \textcolor{Magenta}{\small\textbf{Bloqueado}}}

%% -----------------------------------------------------------------
%% CAJAS tcolorbox
%% -----------------------------------------------------------------
\tcbset{
  cajaNota/.style={
    colback=AzulDigital!10, colframe=AzulDigital,
    fonttitle=\FuenteTitulo\small\bfseries,
    coltitle=white, boxrule=0pt, leftrule=4pt,
    title=Nota, left=8pt, top=6pt, bottom=6pt
  },
  cajaAlerta/.style={
    colback=Magenta!8, colframe=Magenta,
    fonttitle=\FuenteTitulo\small\bfseries,
    coltitle=white, boxrule=0pt, leftrule=4pt,
    title=Advertencia, left=8pt, top=6pt, bottom=6pt
  },
  cajaResultado/.style={
    colback=AzulITM!8, colframe=AzulITM,
    fonttitle=\FuenteTitulo\small\bfseries,
    coltitle=white, boxrule=0pt, leftrule=4pt,
    title=Resultado, left=8pt, top=6pt, bottom=6pt
  },
  cajaDatos/.style={
    colback=GrisClaro, colframe=GrisLinea,
    boxrule=0.6pt, left=8pt, top=6pt, bottom=6pt
  }
}

%% -----------------------------------------------------------------
%% CÓDIGO lstlisting
%% -----------------------------------------------------------------
\lstdefinestyle{coriotcode}{
  basicstyle=\FuenteCodigo\small,
  backgroundcolor=\color{GrisClaro},
  commentstyle=\color{GrisMedio}\itshape,
  keywordstyle=\color{AzulITM}\bfseries,
  stringstyle=\color{Magenta},
  numberstyle=\tiny\color{GrisMedio},
  numbers=left, numbersep=10pt,
  breaklines=true, breakatwhitespace=false,
  frame=none, tabsize=2,
  showstringspaces=false
}
\lstset{style=coriotcode}

%% -----------------------------------------------------------------
%% MACRO: FIRMA SIMPLE
%% -----------------------------------------------------------------
\newcommand{\FirmaSimple}{%
  \vspace{2.5em}
  \noindent
  \begin{tabular}{p{7cm}}
    \hline\\[-6pt]
    {\small\color{GrisPizarra}\DocAutor}\\
    {\small\color{GrisMedio}\DocCargo}\\
    {\small\color{GrisMedio}CORIOTLAB --- ITM}\\
    {\small\color{GrisMedio}\DocFecha}
  \end{tabular}
}

%% =================================================================
%%  INICIO DEL DOCUMENTO
%% =================================================================
\begin{document}

%% -----------------------------------------------------------------
%% === VARIANTE A: PERIÓDICO === descomenta esta sección ===
%% -----------------------------------------------------------------

\PortadaActividades{Informe \TipoPeriodo}{\DocNumero}

\section{Registro de Actividades}

\setlength{\LTpre}{6pt}
\setlength{\LTpost}{6pt}

\begin{longtable}{
  >{\centering\arraybackslash}p{0.55cm}
  p{3.4cm}
  p{5.6cm}
  >{\centering\arraybackslash}p{1.6cm}
  >{\centering\arraybackslash}p{0.9cm}
  >{\centering\arraybackslash}p{2.2cm}
}
\rowcolor{AzulITM}
\color{Blanco}\FuenteTitulo\small\bfseries N\textsuperscript{o} &
\color{Blanco}\FuenteTitulo\small\bfseries Actividad &
\color{Blanco}\FuenteTitulo\small\bfseries Descripción &
\color{Blanco}\FuenteTitulo\small\bfseries Fecha &
\color{Blanco}\FuenteTitulo\small\bfseries Hrs &
\color{Blanco}\FuenteTitulo\small\bfseries Estado \\
\endfirsthead
\rowcolor{AzulITM}
\color{Blanco}\FuenteTitulo\small\bfseries N\textsuperscript{o} &
\color{Blanco}\FuenteTitulo\small\bfseries Actividad &
\color{Blanco}\FuenteTitulo\small\bfseries Descripción &
\color{Blanco}\FuenteTitulo\small\bfseries Fecha &
\color{Blanco}\FuenteTitulo\small\bfseries Hrs &
\color{Blanco}\FuenteTitulo\small\bfseries Estado \\
\endhead
\rowcolor{white}
1 & Configuración sensores BME280 &
Instalación y calibración de los sensores de temperatura, humedad y presión en los nodos ESP32. &
16/06 & 4 & \EstadoCompletado \\
\rowcolor{GrisClaro}
2 & Desarrollo firmware MQTT &
Programación del protocolo MQTT para transmisión de datos al broker HiveMQ local. &
17/06 & 6 & \EstadoCompletado \\
\rowcolor{white}
3 & Integración con Node-RED &
Creación de flujos de procesamiento y visualización de datos en dashboard Node-RED. &
18/06 & 5 & \EstadoEnCurso \\
\rowcolor{GrisClaro}
4 & Pruebas de conectividad WiFi &
Validación del rango de cobertura y estabilidad de conexión en las tres zonas del laboratorio. &
19/06 & 3 & \EstadoEnCurso \\
\rowcolor{white}
5 & Documentación de API REST &
Especificación de endpoints para consulta de datos históricos desde la base InfluxDB. &
20/06 & 2 & \EstadoPendiente \\
\rowcolor{GrisClaro}
6 & Revisión de seguridad TLS &
Implementación de certificados TLS para cifrado del canal MQTT. Pendiente de acceso al servidor. &
--- & --- & \EstadoBloqueado \\
\end{longtable}

\noindent{\small\color{GrisMedio}
\textbf{Total horas registradas:} 20\,h
\quad|\quad
\textbf{Período:} \DocPeriodo
}

\section{Observaciones}

Durante la semana se presentaron dificultades con la estabilidad de la red WiFi del laboratorio en la zona norte. Se identificó interferencia en el canal 6; se realizó cambio al canal 11, con mejora significativa en la latencia (\textasciitilde12\,ms vs.\ 45\,ms previos).

El broker MQTT HiveMQ presentó un reinicio no planificado el 18/06 durante las pruebas de integración, generando pérdida de datos por 20 minutos. Se activó el mecanismo de persistencia de mensajes (QoS 1) como medida correctiva.

\begin{tcolorbox}[cajaNota]
  El módulo de seguridad TLS requiere acceso de administrador al servidor de certificados del ITM. Se solicitó acceso el 19/06 a la Dirección de Sistemas --- tiempo estimado de respuesta: 3 días hábiles.
\end{tcolorbox}

\section{Próximos Pasos}

\begin{itemize}[label={\color{AzulITM}\textbf{>}}, leftmargin=1.8em, itemsep=4pt]
  \item Completar integración Node-RED con visualización en tiempo real de los 8 nodos activos.
  \item Finalizar documentación de la API REST con ejemplos de consulta en Python.
  \item Resolver acceso TLS e implementar cifrado en el canal MQTT.
  \item Realizar prueba de carga con los 8 nodos transmitiendo simultáneamente a 1\,Hz.
  \item Preparar informe de resultados preliminares para reunión del 27/06.
\end{itemize}

\FirmaSimple

%% -----------------------------------------------------------------
%% FIN VARIANTE A
%% -----------------------------------------------------------------


%% -----------------------------------------------------------------
%% === VARIANTE B: GUÍA === descomenta esta sección ===
%% (Comenta todo lo anterior desde \PortadaActividades hasta
%%  \FirmaSimple y descomenta el bloque siguiente)
%% -----------------------------------------------------------------

%\PortadaActividades{Guía Técnica}{\DocNumero}
%
%\section{Objetivo de la Guía}
%
%Esta guía describe el proceso de configuración e integración de sensores BME280
%con microcontroladores ESP32 usando MicroPython y el protocolo MQTT para la
%transmisión de datos hacia un broker local en la red del CORIOTLAB.
%
%\section{Requisitos Previos}
%
%\begin{itemize}[label={\color{AzulITM}\textbf{>}}, leftmargin=1.8em, itemsep=4pt]
%  \item Microcontrolador ESP32 con MicroPython instalado (v1.22 o superior).
%  \item Sensor BME280 conectado por I\textsuperscript{2}C (SDA=GPIO21, SCL=GPIO22).
%  \item Broker MQTT activo en la red local (HiveMQ CE o Mosquitto).
%  \item Python 3.10+ con las librerías \texttt{paho-mqtt} y \texttt{bme280}.
%  \item Acceso a la red WiFi del laboratorio (SSID: CORIOTLAB-LAB).
%\end{itemize}
%
%\section{Pasos}
%
%\subsection{Instalación del firmware MicroPython en ESP32}
%
%Descargar el firmware desde el sitio oficial y flashear con \texttt{esptool.py}:
%
%\begin{tcolorbox}[cajaDatos]
%\begin{lstlisting}[language=bash]
%# Borrar flash previa
%esptool.py --port /dev/ttyUSB0 erase_flash
%# Flashear firmware MicroPython
%esptool.py --chip esp32 --port /dev/ttyUSB0 \
%  write_flash -z 0x1000 esp32-micropython-1.22.0.bin
%\end{lstlisting}
%\end{tcolorbox}
%
%\subsection{Lectura del sensor BME280}
%
%Copiar el siguiente script al ESP32 usando \texttt{ampy} o Thonny IDE:
%
%\begin{tcolorbox}[cajaDatos]
%\begin{lstlisting}[language=Python]
%import bme280
%from machine import I2C, Pin
%
%i2c = I2C(0, sda=Pin(21), scl=Pin(22), freq=400000)
%sensor = bme280.BME280(i2c=i2c)
%
%temp, pres, hum = sensor.read_compensated_data()
%print(f"Temp: {temp/100:.1f} C | Hum: {hum/1024:.1f}% | Pres: {pres/25600:.1f} hPa")
%\end{lstlisting}
%\end{tcolorbox}
%
%\begin{tcolorbox}[cajaNota]
%  Los valores del BME280 vienen en formato de punto fijo.
%  Temperatura / 100, humedad / 1024 y presión / 25600 para obtener unidades SI.
%\end{tcolorbox}
%
%\subsection{Publicación de datos por MQTT}
%
%\begin{tcolorbox}[cajaDatos]
%\begin{lstlisting}[language=Python]
%from umqtt.simple import MQTTClient
%import json, time
%
%client = MQTTClient("esp32_node1", "192.168.1.100")
%client.connect()
%
%while True:
%    payload = json.dumps({"temp": temp/100, "hum": hum/1024})
%    client.publish(b"coriot/nodo1/ambiental", payload.encode())
%    time.sleep(5)
%\end{lstlisting}
%\end{tcolorbox}
%
%\begin{tcolorbox}[cajaAlerta]
%  No usar QoS 0 en producción. Configurar QoS 1 para garantizar
%  la entrega de mensajes ante pérdidas momentáneas de conexión.
%\end{tcolorbox}
%
%\section{Conclusión}
%
%Con estos pasos el nodo ESP32 queda operativo y transmitiendo datos ambientales
%cada 5 segundos al broker MQTT. El siguiente paso es configurar Node-RED para
%suscribirse al tópico \texttt{coriot/\#} y visualizar los datos en tiempo real.
%
%\FirmaSimple

%% -----------------------------------------------------------------
%% FIN VARIANTE B
%% -----------------------------------------------------------------

\end{document}
